%==============================================================================
%  §4  Bridging Tree Topology to Linear Algebra                    [owner: A1]
%  Source: v8 §4 (sec:bridging, v8:344-471).
%  Moved out:  def:imbalance + eq:imbalance_ratio  -> sec:problem (inline)
%              asm:margin    + eq:margin_assumption -> sec:problem
%
%  [Reviewer Note (v9 refine pass): section retitled "Topological Scaling and
%  Spectral Gap" -- the old title named the activity rather than the content,
%  and "bridging X to Y" reads as a lecture heading.  The three \subsection*
%  headings ("Imbalance raises...", "Imbalance erodes...", "Two failure modes")
%  are removed: each announced in a sentence what the proposition or lemma
%  immediately below it states formally, so the reader was told everything
%  twice.  Proposition 1 and Lemma 1 are now stated directly, and the two
%  failure modes run as continuous prose rather than under \paragraph headings.]
%==============================================================================
\section{Topological Scaling and Spectral Gap}\label{sec:bridging}

The reader now has the population model and the recovery problem. What remains is
to explain \emph{why} the topology of the primary split governs recovery. We
answer this through the single structural parameter $\eta$ of
\eqref{eq:imbalance_ratio} and trace its effect on the two spectral quantities
the main theorem depends on. Working with the clean population Laplacian
$L=D-S$, both the subspace coherence $\coh$ and the spectral gap
$\Dlam=\lambda_3-\lambda_2$ are fixed by macro-topological features alone, and
$\eta$ moves both against recovery: it first inflates $\coh$, then erodes
$\Dlam$.

\begin{proposition}[Topological scaling of coherence]\label{prop:coherence}
Under \Cref{asm:cfbm}, the tree imbalance $\eta$ maps deterministically to the
subspace coherence,
\begin{equation}\label{eq:coherence_linear_growth}
   \coh=\frac{1+\eta}{2}.
\end{equation}
\end{proposition}
\noindent The identity follows by substituting the Fiedler coordinates of
\Cref{lem:A-decomp} into \eqref{eq:coherence_definition}; the short computation is given in
\Cref{app:comp1}. Thus a balanced split attains the
optimal coherence $\coh=1$, and the signal spreads more unevenly across nodes as
$\eta$ grows; the generated regime of \Cref{sec:empirical} confirms the identity
on data from the full generative pipeline. The identity is reported as a
structural diagnostic: since the noise bound of \Cref{lem:B-bernstein} is obtained
from the row norms of $S$ directly (\Cref{app:comp2}), $\coh$ does not itself enter
the sampling rate, and the $\eta$-dependence of \eqref{eq:main_rate} arrives
through the margin and the gap instead.

% [Reviewer Note: the coherence-vs-eta validation figure (former Figure 2,
%  coherence_vs_eta_by_n.png) was removed at the author's instruction.  The
%  sentence that carried it now points forward to the generated regime in §6,
%  where the same identity is reported as a measurement.]

The second spectral quantity is the gap $\Dlam=\lambda_3-\lambda_2$ that isolates
the Fiedler vector, and the same imbalance that raises $\coh$ drives it toward
zero.

\begin{lemma}[Spectral-gap lower bound]\label{lem:gap}
Let $L$ be the combinatorial Laplacian of a similarity matrix satisfying the
molecular clock, with structural margin $\rho=\Sin^{\min}-\Sout^{\max}$. For
$m_{\min}=\min(a,b)$ and imbalance $\eta$, the spectral gap isolating the Fiedler
vector satisfies
\begin{equation}\label{eq:spectral_gap_margin}
   \Dlam \ge m_{\min}\rho-\eta\,m_{\min}\,\Sout^{\max}.
\end{equation}
\end{lemma}
\noindent The full algebraic-connectivity derivation is given in \Cref{app:comp1},
building on the eigenpairs of \Cref{lem:A-decomp}. The bound
\eqref{eq:spectral_gap_margin} is purely a CFBM statement, controlling the gap
through the abstract minimum intra-clan similarity $\Sin^{\min}$ of
\eqref{eq:cfbm_floor}. Under the sufficient condition of
\Cref{asm:hbm} that minimum stops being a free parameter and becomes
a function of tree shape: it is pinned to the geometric floor
$\Sin\,\alpha^{D_{\max}}$, so depth erodes the gap in the same way imbalance does,
and the gap survives only while the decay $\alpha^{D_{\max}}$ still clears the
imbalance-weighted boundary --- quantified in \Cref{prop:D-gap}.

\Cref{prop:coherence} and \Cref{lem:gap} show that the same parameter $\eta$ acts
on both spectral quantities at once, which produces two distinct ways for
top-down recovery to fail; together they bound the admissible region of
$(m,\eta)$. The first is \emph{structural collapse},
$\Dlam\to0$, which we state here as a corollary of the gap lemma; the second is
\emph{statistical infeasibility}, $p>1$, which is a consequence of
\Cref{thm:main-sim} below. A positive gap requires the structural margin to
dominate the imbalance-weighted cross-clan boundary,
$m_{\min}\rho>\eta\,m_{\min}\Sout^{\max}$, i.e.\
$\eta<\Sin^{\min}/\Sout^{\max}$ --- precisely what \Cref{asm:margin} records.
Translating this
ratio into tree distances bounds the imbalance the topology can sustain.

\begin{corollary}[Structural imbalance tolerance]\label{cor:tolerance}
Write similarities as distances, $S_{ij}=e^{-d(x_i,x_j)}$, and let
$r(\mathcal T)=\max_{i,j}(-\log S_{ij})$ be the diameter and $h(\mathcal T)$ the
depth of $\mathcal T$ \cite{aizenbud2023spectral}, with $h(\mathcal T)<r(\mathcal T)$.
Then a positive spectral gap is possible only if
\begin{equation}\label{eq:eta_distance_bound}
   \eta\le\mathcal O\!\bigl(e^{\mathcal O(h(\mathcal T))-r(\mathcal T)}\bigr).
\end{equation}
In particular, for an average Kingman coalescent tree, where
$r(\mathcal T),h(\mathcal T)=\mathcal O(\log m)$
\cite{erdos1999logs, aldous2001stochastic}, this gives $\eta\le\mathcal O(m^c)$ for
a model-dependent constant $c$.
\end{corollary}
\noindent The proof sketch (substituting the diameter/depth bounds into
$\eta<\Sin^{\min}/\Sout^{\max}$) is given in \Cref{app:comp1}.

% [Reviewer Note (2026-08-04): this paragraph previously re-explained the
%  statistical failure mode -- localisation of the signal, the coherence-limited
%  matrix-completion barrier, the same three citations -- in wording near
%  identical to rem:infeasible-mech in sec:main-result-dist.  Two faults, not
%  one: the two-mode framing is already stated above (before cor:tolerance) and
%  in the introduction, and the mechanism was narrated here before thm:main-sim
%  and cor:infeasible had been stated at all.  The explanation now occurs once,
%  at the bound that exhibits it; this paragraph keeps only the forward pointer.]
So under balanced generative models the gap survives polynomial imbalance. A
positive gap is necessary but not sufficient: even when $\Dlam>0$, the rate
demanded by \Cref{thm:main-sim} grows with $\eta$, and \Cref{cor:infeasible}
locates the imbalance past which that demand exceeds $p=1$.
